\documentclass[10pt,letterpaper]{article}
\usepackage[margin=0.68in,headheight=14pt,headsep=12pt,footskip=24pt]{geometry}
\usepackage[T1]{fontenc}
\usepackage{lmodern,amsmath,booktabs,tabularx,graphicx,xcolor,fancyhdr,enumitem}
\usepackage{pgfplots}
\usepackage[colorlinks=true,urlcolor=teal!65!black,linkcolor=teal!65!black,citecolor=teal!65!black]{hyperref}
\pgfplotsset{compat=1.18}
\definecolor{ink}{HTML}{183244}
\definecolor{accent}{HTML}{087F8C}
\definecolor{pale}{HTML}{EFF6F7}
\setlength{\parindent}{0pt}
\setlength{\parskip}{5pt}
\setlength{\tabcolsep}{4pt}
\renewcommand{\arraystretch}{1.12}
\setlist[itemize]{leftmargin=13pt,itemsep=1pt,topsep=2pt}
\pagestyle{fancy}
\fancyhf{}
\fancyhead[L]{\small\sffamily\color{ink}CancerOmicsLake}
\fancyhead[R]{\small\sffamily Research brief \textbar\ 19 October 2026}
\fancyfoot[L]{\small Pronab Chandra Roy}
\fancyfoot[R]{\small\thepage\ / 4}
\renewcommand{\headrulewidth}{0.3pt}
\newcommand{\topic}[1]{\par\vspace{5pt}{\large\sffamily\bfseries\color{ink}#1}\par\vspace{1pt}}
\newcommand{\subtopic}[1]{\par\vspace{2pt}{\bfseries\color{ink}#1}\par}
\newcommand{\takeaway}[1]{\par\smallskip\noindent\colorbox{pale}{\parbox{\dimexpr\linewidth-2\fboxsep\relax}{\textbf{Finding.} #1}}\par}
% Generated from checked aggregate evidence.
\newcommand{\TCGARows}{36,456,660}
\newcommand{\GTExRows}{11,240,000}
\newcommand{\NativeScripts}{7}
\newcommand{\WorkflowCells}{21}
\newcommand{\MatchedTCGA}{417}
\newcommand{\PrimaryNormals}{150}
\newcommand{\SharedGenes}{55,660}
\newcommand{\MatchedTotal}{617}
\newcommand{\MatchedGTEx}{200}
\newcommand{\SigmoidNormals}{50}
\newcommand{\MatchedRows}{BRCA & 162 & 50 & 0.559 & 89.85\% & 23,247 & 0.0046 \\
COAD & 119 & 50 & 0.556 & 90.53\% & 22,007 & 0.0026 \\
LUAD & 136 & 50 & 0.571 & 90.21\% & 23,732 & 0.0058 \\}
\newcommand{\TieMin}{412}
\newcommand{\TieMax}{515}
\newcommand{\NativeFailed}{63}
\newcommand{\PairedFailed}{60}
\newcommand{\ExactComparisons}{46}
\newcommand{\ExactAttempts}{6,624}
\newcommand{\ExactFailed}{36}
\newcommand{\ContractCases}{32}
\newcommand{\ContractDecisions}{192}
\newcommand{\RecomputeEvents}{400}
\newcommand{\RecomputeWorkloads}{40}
\newcommand{\CertifiedEvents}{233}
\newcommand{\LifecycleWins}{0}
\newcommand{\AbstentionDatasets}{3,600}
\newcommand{\JointWrong}{47}
\newcommand{\HolmWrong}{5}
\newcommand{\FragilityWorlds}{112}
\newcommand{\FragilityCases}{7}
\newcommand{\TieDraws}{1,000}
\newcommand{\BioBootstraps}{200}
\newcommand{\BootstrapTieDraws}{100}
\newcommand{\DirectionThreshold}{0.2}
\newcommand{\DetectionPercent}{20}

\hypersetup{pdftitle={CancerOmicsLake: engineering and sensitivity of cancer-gene rankings},pdfauthor={Pronab Chandra Roy}}
\begin{document}

{\sffamily\color{accent}\small BIOINFORMATICS DATA ENGINEERING \ / \ PHD INTERVIEW SUMMARY}\par
{\sffamily\bfseries\color{ink}\LARGE CancerOmicsLake}\par
{\sffamily\color{ink}\Large An auditable platform for studying the sensitivity\newline of cancer-gene rankings}\par
\textbf{Pronab Chandra Roy}\hfill\href{https://github.com/pronabroy360/CancerOmicsLake}{Code and project repository}\par

\topic{1. Research question and engineering contribution}
\textbf{How much can a cancer-gene ranking change when its reference population,
expression processing, tie handling, or numerical environment changes, even when
the workflow executes reproducibly?} CancerOmicsLake addresses this question by
connecting public data acquisition, harmonization, analytical tables and recorded
experiments. Its strongest contribution is the implementation and evaluation of
an inspectable research workflow, including evidence of its limitations.

\subtopic{Data and scope}
The project integrates open-access GDC/TCGA metadata, gene expression, clinical
and biospecimen fields, somatic mutation summaries, GTEx normal-tissue expression,
and publicly released recount3 processed expression~\cite{recount}. The cancer
cohorts are breast invasive carcinoma (BRCA), lung adenocarcinoma (LUAD) and colon
adenocarcinoma (COAD). Normal references are breast, lung, transverse colon and
sigmoid colon. The frozen primary expression inventory contains
\textbf{\TCGARows\ TCGA} and \textbf{\GTExRows\ GTEx} measurement rows;
these are gene--specimen measurements, not counts of patients or sequencing files.

\begin{center}
\fcolorbox{accent}{pale}{\parbox{0.94\linewidth}{\centering\small
\textbf{GDC / GTEx / recount3 / annotations}\\[3pt]
$\downarrow$\quad configured acquisition, retries, checksums and source metadata\\[3pt]
\textbf{Bronze: immutable inputs}\ $\longrightarrow$\
\textbf{Silver: typed, harmonized Parquet}\ $\longrightarrow$\
\textbf{Gold: analytical marts}\\[3pt]
$\downarrow$\quad DuckDB / dbt tests / quality gates / frozen experiment manifests\\[3pt]
\textbf{SQL, API, dashboard, graph exports and claim-to-source evidence}
}}\end{center}

\subtopic{Implemented capabilities}
Python/Polars ingestion and transformation support configurable cohorts and
bounded downloads; DuckDB and dbt provide queryable dimensions and facts.
Orchestration, Docker and CI support repeatable runs. Streamlit, FastAPI and
Neo4j/Graphify exports expose the results. Validation checks cover access,
checksums, types, duplicate keys, gene/sample joins and relationships.

The analytics include expression and reference comparisons, consequence-stratified
mutation frequencies, candidate prioritization, pathway overrepresentation with
Benjamini--Hochberg correction within cancer/candidate set, and graph degree
summaries. These provide exploratory context; pathway membership, mutation
frequency and graph connectivity do not by themselves validate a candidate.

\subtopic{Research process and contribution}
Work progressed from infrastructure to reference sensitivity, exact specimen
matching, tie diagnostics, published baselines and numerical replay. Repairs
addressed duplicate measurements, Ensembl version normalization while preserving
\texttt{\_PAR\_Y} loci, and donor/case leakage across partitions. Plans and input
hashes preserve earlier evidence when a later study changes the design.
Individual manifests and per-gene intermediates stay outside the aggregate brief.

\takeaway{A functioning data platform makes the assumptions and failure modes of
candidate rankings inspectable. The studies below assess what that infrastructure
can establish scientifically.}

\newpage
\topic{2. Matched specimens: broad agreement, unstable leading lists}
\subtopic{Study design}
Matched-processing v1 validated \textbf{\MatchedTotal\ exact specimens}:
\MatchedTCGA\ TCGA tumors and \MatchedGTEx\ GTEx normals. Matching checked GDC
sample identity or GTEx SAMPID, case/donor, tissue, QC and count-column mapping.
Ambiguities were excluded and corrected measurement keys were required to be
unique. The three primary comparisons use \PrimaryNormals\ normals; the remaining
\SigmoidNormals\ sigmoid-colon specimens are retained separately. COAD uses transverse colon
for the primary comparison. Both processing pipelines use identical specimens.

Native TPM and recount3 AUC-scaled coverage are kept in their respective units.
V1 compared within-specimen gene ranks, detection and top-expressed genes, then
compared tumor--normal effects. V2, designed after inspecting v1, examined ties,
detection, annotation and workflow provenance. It is an exploratory follow-up.

For each gene and pipeline, the signed rank-biserial effect is
\[
 r_g=\frac{1}{n_Tn_N}\sum_{i=1}^{n_T}\sum_{j=1}^{n_N}
 \operatorname{sign}(x_{gi}^{T}-x_{gj}^{N}).
\]
This measures relative ordering within a pipeline. It does not make expression
units comparable or provide a differential-expression significance test.

\begin{center}\small
\textbf{Matched v2: complete shared-gene universe (\SharedGenes\ genes per cancer)}\\[4pt]
\begin{tabular}{lrrrcrr}
\toprule
Cancer & Tumors & GTEx & $\tau_b$ & Direction agreement & Eligible genes & Mean $J_{50}$\\
\midrule
\MatchedRows
\bottomrule
\end{tabular}\end{center}
{\footnotesize $\tau_b$: signed-effect Kendall correlation. Direction agreement
uses genes with $|r_g|\geq\DirectionThreshold$ in both pipelines; its denominator is shown.
$J_{50}$: Jaccard overlap of absolute-effect top-50 lists after seeded boundary-tie
resolution. Specimens are not the denominator for these gene-level metrics.}\par

\begin{center}
\begin{tikzpicture}
\begin{axis}[width=0.46\linewidth,height=1.50in,ybar,bar width=14pt,
 ymin=0,ymax=1,ytick={0,0.5,1},symbolic x coords={BRCA,COAD,LUAD},xtick=data,
 title={\small Broad signed-effect agreement},ylabel={\small Kendall $\tau_b$},
 tick label style={font=\footnotesize},title style={font=\small},
 axis lines*=left,grid=major,grid style={gray!15}]
\addplot[fill=accent,draw=accent] table[x=cancer,y=tau]{matched_metrics.dat};
\end{axis}
\end{tikzpicture}\hfill
\begin{tikzpicture}
\begin{axis}[width=0.46\linewidth,height=1.50in,ybar,bar width=14pt,
 ymin=0,ymax=0.01,ytick={0,0.005,0.01},scaled y ticks=false,
 yticklabel style={/pgf/number format/fixed,/pgf/number format/precision=3},
 symbolic x coords={BRCA,COAD,LUAD},xtick=data,
 title={\small Leading-list overlap},ylabel={\small Mean top-50 Jaccard},
 tick label style={font=\footnotesize},axis lines*=left,grid=major,grid style={gray!15}]
\addplot[fill=ink,draw=ink] table[x=cancer,y=jaccard]{matched_metrics.dat};
\end{axis}
\end{tikzpicture}
\end{center}
{\footnotesize Figure: the panels use different vertical scales and different
statistics; their bar heights should not be compared directly. Values come from
the same matched comparisons in the table.}\par

V2 evaluated all shared genes and a fixed informative subset detected ($>0$) in
at least \DetectionPercent\% of both groups in both pipelines. It used \TieDraws\ seeded tie resolutions
and \BioBootstraps\ paired case/donor bootstrap replicates, with \BootstrapTieDraws\ tie resolutions within
each replicate. Tie-break uncertainty and descriptive sampling intervals were
reported separately. Filtering improved broad concordance but left leading-list
overlap weak. Native boundary scores tied across \TieMin--\TieMax\ genes.

\takeaway{Genome-wide agreement does not establish stable short candidate lists.
The comparison measures processing sensitivity conditional on available matched
specimens; it does not isolate a causal source of the discrepancy.}

\newpage
\topic{3. Reference choice, numerical execution and external baselines}
\subtopic{Ranking candidates against two normal references}
Globally separated case/donor partitions compared adjacent-normal and GTEx
references in corrected native and paired study designs. A bootstrap
lower-quantile score was evaluated against single-reference effects, mean ranks,
reference maximin and the official R \texttt{RobustRankAggreg} implementation~\cite{rra}.
Top-25/50/100 lists were evaluated with heldout effect/direction summaries and
paired resampling. Bootstrap rank stability, reference ablations and reference
discordance were retained as exploratory diagnostics.

The proposed score showed \textbf{no top-50 advantage over simple maximin} in
the corrected studies. RRA supplies an established ranking baseline; because
both input lists share tumor specimens, its independent-list null cannot justify
calibrated significance here. A constructed shared-normal-bias example further
showed that agreement between references can coexist with incorrect biological
selection. These results restrict the interpretation of consensus scores.

\subtopic{Execution can reproduce candidates while failing numerical acceptance}
Archived replay compared macOS/Linux with pinned dependencies and Polars
1.42.1/1.44.2. It recorded ordering, tie counts and intermediate arrays, while
retaining the original tolerances: $\mathrm{rtol}=10^{-8}$,
$\mathrm{atol}=10^{-10}$. The three historical failed rows repeated one Spearman
statistic at three top-k values. The earliest observed platform divergence was
in log2 effect arrays, followed by tie-sensitive ranks; the underlying platform
cause was not established by the dependency comparison.

Corrected global studies reproduced candidate lists and heldout metrics, but
\textbf{\NativeFailed\ native and \PairedFailed\ paired sensitivity cells failed}
the fixed tolerance. A later exact-arithmetic study compared existing eager
rational, specialized rational and CGAL lazy exact implementations: all
\ExactComparisons\ workload/specification comparisons agreed, including native
macOS checks. Of \ExactAttempts\ process attempts, \ExactFailed\ binary64 attempts
failed on a constructed extreme-value fixture; exact-method attempts succeeded.
This separate study did not repair the archived numerical failures or establish
a new arithmetic algorithm.

\subtopic{Fair external workflow and failure experiments}
The pinned \texttt{tcga-data-nf} preparation implementation~\cite{workflow} was
compared on cached recount3 inputs with matching specimen/gene membership and
normalization/filtering semantics. There were \WorkflowCells\ comparable
preparation workloads and \NativeScripts\ successful native-script integrations.
Output equivalence was established before timing. A common R normalization
backend, adapter costs and emulated hardware limit performance interpretation;
no end-to-end Nextflow or engine-speed advantage was established.

\begin{center}\small
\begin{tabularx}{\linewidth}{@{}>{\raggedright\arraybackslash}p{0.25\linewidth}>{\raggedright\arraybackslash}X>{\raggedright\arraybackslash}X@{}}
\toprule
Tested safeguard & Observed behavior & Interpretation\\
\midrule
Corrupt input, missing metadata, conflicting duplicates & Benchmark adapter/parser rejected each & Guards worked at the recorded validation layer.\\
Interrupted execution & No incorrect output committed; fresh retry recovered & Recovery tested with an explicit pause hook.\\
Stale cached evidence & Evidence hash guard rejected reuse & External native stale-cache handling was unmeasured.\\
Frozen sources and reports & Hashes, manifests and earlier failures retained & Traceability supports audit; it does not prove biological truth.\\
\bottomrule
\end{tabularx}\end{center}

\takeaway{The supported contribution is a measured engineering case study:
source-to-claim traceability, controlled comparisons and preservation of failed
scientific checks. Passing software tests does not override those failures.}

\newpage
\topic{4. Additional method experiments and research status}
The project also tested whether its failure cases motivated a sharper method or
systems contribution. These exploratory studies used frozen evidence and strong
baselines. Their bounded negative results are part of the research record.

\begin{center}\small
\begin{tabularx}{\linewidth}{@{}>{\raggedright\arraybackslash}p{0.23\linewidth}>{\raggedright\arraybackslash}X>{\raggedright\arraybackslash}X@{}}
\toprule
Study & Executed evidence & Outcome and boundary\\
\midrule
Executable evidence contracts & \ContractCases\ synthetic cases; \ContractDecisions\ tool/case/pass decisions using SQL, SHACL and Snakemake & All matched the prototype. Each admitted two false attestations; metadata consistency cannot establish original truth.\\
Decision-aware recomputation & \RecomputeEvents\ distinct patch events across \RecomputeWorkloads\ workloads; \CertifiedEvents\ unchanged sets certified & No decision mismatches; total cost beat exact incremental updating in \LifecycleWins/\RecomputeWorkloads\ workloads.\\
Exact ranking arithmetic & \ExactComparisons\ comparisons using established exact methods & Exact results agreed; no unique capability or uniformly fastest implementation.\\
Rank selection with abstention & \AbstentionDatasets\ synthetic datasets; Holm, Bonferroni and joint bootstrap & Narrow label screen passed; invalid pair directions occurred in \JointWrong/200 versus \HolmWrong/200 datasets in one moderate-signal cell.\\
Minimum claim-flipping changes & \FragilityWorlds\ synthetic workflow alternatives in \FragilityCases\ scenarios & SQL and Python agreed on all decisions and witnesses; no unique capability demonstrated.\\
\bottomrule
\end{tabularx}\end{center}
{\footnotesize Repeated timings/tool passes are not independent datasets.
The abstention diagnostic is the 30-donor moderate-signal cell. Its pairwise
error undermines a calibrated-confidence claim despite correct observed top-k
labels. The official confident-feature-ranking implementation~\cite{confident}
was not benchmarked. The claim-flipping test replayed prepared synthetic
alternatives rather than generating real quantification workflows.}\par

\subtopic{Position relative to prior work}
Arora et al.~\cite{arora} already established that processing differences affect
expression and downstream results. CancerOmicsLake adds an auditable combination
of matched candidate-effect comparisons, tie diagnostics and explicit numerical
failure records; it does not claim to discover processing sensitivity first.
Direct Arora replication remains unexecuted: the checked expression-archive
locations returned access errors on 16 October 2026. This does not establish
that the matrices were never published. Their endpoints and ours differ.

\subtopic{What this research supports, and the next step}
The deliverable is a working cancer-omics data platform and a documented
engineering/sensitivity case study. It has not established a superior new
ranking method, independent biological validation or clinical utility. Batch
effects, available-sample selection and annotation differences remain limitations.
Existing inspected holdouts cannot provide fresh confirmation.

The next concrete steps are independent investigator reproduction of the
prepared handoff, external comparison on identical endpoints, and an untouched
evaluation cohort selected before method development. Any future algorithm
needs a clear estimand and benefit over maximin or competent systems baselines.
For a PhD, my interest is developing and evaluating reliable data infrastructure
for biomedical research, with this platform as an implemented starting point.

{\footnotesize\textbf{Document status.} Working research summary for interview
discussion, updated from the frozen v5 evidence and subsequent feasibility
studies. Code and drafting used AI assistance. Author review, affiliation,
correspondence and formal declarations remain to be finalized; independent
human reproduction is pending. The evidence ledger and reproducible LaTeX
build accompany this brief.}\par

\vspace{-3pt}
\begin{thebibliography}{5}\footnotesize\raggedright\setlength{\itemsep}{0pt}\setlength{\parskip}{0pt}
\bibitem{recount} Wilks et al. \emph{recount3: summaries and queries for large-scale RNA-seq expression and splicing}. Genome Biology, 2021. \href{https://doi.org/10.1186/s13059-021-02533-6}{doi:10.1186/s13059-021-02533-6}.
\bibitem{rra} Kolde et al. \emph{Robust rank aggregation for gene list integration and meta-analysis}. Bioinformatics, 2012. \href{https://doi.org/10.1093/bioinformatics/btr709}{doi:10.1093/bioinformatics/btr709}.
\bibitem{workflow} Quackenbush Lab. \emph{tcga-data-nf}. \href{https://github.com/QuackenbushLab/tcga-data-nf/tree/1069eb9c01b97640fd532ffac8551096d2a2790a}{Pinned workflow source}; benchmark scope described on p.~3.
\bibitem{confident} Neuhof and Benjamini. \emph{Confident Feature Ranking}. AISTATS, 2024. \href{https://proceedings.mlr.press/v238/neuhof24a.html}{PMLR 238:1468--1476}.
\bibitem{arora} Arora et al. \emph{Variability in estimated gene expression among commonly used RNA-seq pipelines}. Scientific Reports, 2020. \href{https://doi.org/10.1038/s41598-020-59516-z}{doi:10.1038/s41598-020-59516-z}.
\end{thebibliography}
\end{document}
